\section{Participation, Termination, and Portfolio Permissions}
\label{sec:r9_contract}
The operating obligation and the investment permissions are distinct economic primitives. An agent who can leave a costly activity need not value an adjustment technology as does an agent committed to continue. Likewise, a change in the permitted size of a position need not have the same effect as a change in preferences. We embed the original operating economy in a participation contract, characterize the price of relaxing its obligation, and identify the investment-permission margins that move the two indifference loci. The original noncancellable experiment remains a member of this family.

\subsection{An Operating Mandate with Participation and Limited Enforcement}
\label{subsec:r9_participation}
At date zero, an agent may accept an operating mandate or take the outside payoff $G(x_0)$. Acceptance fixes the managed initial state $x_0=(2,1.25)$, the rebalancing calendar, and the stochastic settlement technology of Section \ref{subsec:finite_protocol}. A separate, additively separable utility numeraire supports a signing grant $q$ and an enforceable termination charge $F$. Neither payment is credited to managed wealth $X$. Crediting the grant to that state would change consumption, transitions, and the experiment, and is not the contract considered here.

The mandate stipulates a minimum operating term $m h$, where $h=1/8$ and $m\in\{1,\ldots,8\}$. At an interior rebalancing date $n\geq m$, the agent may surrender for $G(x)-F$. Before that date, operation is compulsory. At a resource or preference boundary the mandate is discharged without a charge: the boundary records a contractually verifiable inability to continue within the permissible state domain. This discharge provision differs from an elective surrender at an ordinary interior node. It explains why the same underlying liquidation payoff does not imply the same right to terminate. The case $m=8$ is the original noncancellable mandate, independently of $F$.

The charge is supported by externally enforceable guarantee capacity $F\leq E\leq\bar E$, where $\bar E$ is the available capacity limit. Let $C(E)\geq0$ be the agent's nondecreasing ex ante cost of maintaining that capacity, measured in the separate numeraire. It may represent a surety premium or a collateral carrying cost. It is independent of the within-mandate control conditional on the contracted capacity, and enters the participation payment explicitly. No costless, unbounded borrowing against managed wealth is assumed. A fully collateralized interpretation is also possible if the collateral account and its return are specified separately; its carrying cost cannot be silently omitted. The calculations below first report $C(E)=0$ so that the effect of the operating technology is visible. A specified positive $C(E)$ adds to the required grant whenever participation binds.

These are assumptions of a theoretical procurement economy, not claims about a particular intermediary. They provide a contracting interpretation and a transparent participation test. They do not establish that this mandate, its settlement lottery, or its enforcement technology is optimal among all contracts. In particular, a principal must be willing to pay for the service; declaring the operation compulsory does not itself satisfy participation.

Write $r\in\{\mathrm{adj},0\}$ for the two adjustment regimes and $\sigma\in\{+,-\}$ for the initial position class. The regime $0$ prohibits deliberate adjustment, not stochastic movements of preferences. Let $W^r_\sigma(\lambda,d,F,m)$ denote the operating value before the signing grant and the capacity cost. A deterministic Markov policy $p$ specifies both operating actions and an elective surrender rule. Let $\tau$ be its first discharge, surrender, or terminal date. Write $\mathcal E_p$ for elective surrender before discharge or maturity, and define
\begin{align}
 \mathcal A_p&=\E_p\left[\int_0^\tau e^{-\rho t}\,dt\right],\qquad
 \mathcal S_p=\E_p\left[e^{-\rho\tau}\mathbf 1_{\mathcal E_p}\right],\label{eq:r9_features}\\
 J_p(\lambda,d,F,m)&=\mathcal B_p(\lambda,m)+d\mathcal A_p(\lambda,m)-F\mathcal S_p(\lambda,m).
 \label{eq:r9_stop_affine}
\end{align}
Here $\mathcal B_p$ contains consumption utility, preference-state utility, adjustment cost, and liquidation without the charge. In the finite protocol, the integral in \eqref{eq:r9_features} uses the actual within-period exit annuity. The surrender exposure is discounted incidence, not an undiscounted exit probability. Fixing a policy fixes these features; reoptimizing a contract need not do so.

\begin{proposition}[Participation, service procurement, and the surrender envelope]
\label{prop:r9_participation}
For fixed $(\lambda,m)$ and a fixed feasible action technology, $W^r_\sigma=\max_{p\in\mathcal P^r_\sigma(m)}J_p$ is convex in $(d,F)$, nondecreasing in $d$, nonincreasing and $1$-Lipschitz in $F$. The smallest nonnegative signing grant satisfying participation is
\begin{equation}
 q_r^*(E)=\left[G(x_0)-\max_{\sigma}W^r_\sigma+C(E)\right]_+.
 \label{eq:r9_participation}
\end{equation}
Suppose a principal obtains flow $b$ per unit of discounted operating duration and pays the grant, while termination charges are remitted to a separate enforcement institution. For a specified optimal agent policy $p_r^*$ with $\mathcal A_{p_r^*}>0$, that contract is individually rational for both parties at the minimum grant if and only if
\begin{equation}
 b\geq \frac{q_r^*(E)}{\mathcal A_{p_r^*}}.
 \label{eq:r9_procurement}
\end{equation}
At a point with a unique exposed feature vector, the class-value derivatives are $W_d=\mathcal A_{p^*}$ and $W_F=-\mathcal S_{p^*}$. At a tie, the directional derivative is the maximum of $v_d\mathcal A_p-v_F\mathcal S_p$ over optimizing policies. These statements concern the specified contract, not an optimization over all mandates.
\end{proposition}

The proof is an affine-policy envelope argument, given in Supplement S.13. Equation \eqref{eq:r9_procurement} imposes a substantive participation restriction on the principal as well as on the agent. It compares compensation with delivered service rather than interpreting a low grant for an immediately abandoned activity as a procurement improvement. The displayed values are utility-numeraire quantities. They are not monetary willingness-to-pay estimates for an investor whose entire wealth follows CRRA utility.

Define $\Delta^r=W^r_+-W^r_-$ and $\mathcal O=\Delta^{\mathrm{adj}}-\Delta^0$. Where the feature derivatives exist,
\begin{equation}
 \Delta^r_d=\mathcal A^r_+-\mathcal A^r_-,\qquad
 \Delta^r_F=-\bigl(\mathcal S^r_+-\mathcal S^r_-\bigr).
 \label{eq:r9_delta_derivatives}
\end{equation}
Thus increasing a surrender charge need not increase either the relative adjustment option or the positive-position advantage. The sign depends on which class uses surrender more. If all class-specific participation constraints bind, then
\begin{equation}
 \mathcal O=(q^0_+-q^{\mathrm{adj}}_+)-(q^0_--q^{\mathrm{adj}}_-),
 \label{eq:r9_option_price}
\end{equation}
for a common capacity cost and outside payoff. The relative option is the difference between the compensation savings from adjustment in the two position classes. Adjustment weakly lowers each class's required compensation; it need not lower it more for the positive class.

\subsection{When Does a Finite Charge Implement Continued Operation?}
\label{subsec:r9_enforcement}
Let $V^r_n$ be the original noncancellable continuation value. For the new mandate the recursion at eligible interior nodes is
\begin{equation}
 W^r_n(x)=\max\left\{G(x)-F,\ \max_{a\in A^r_n(x)}
          \bigl[r^{\lambda,d}_n(x,a)+K^\lambda_n(x,a)W^r_{n+1}\bigr]\right\}.
 \label{eq:r9_stopping}
\end{equation}
Before date $m$ only the second entry is available. Boundary nodes continue to pay $G$, terminal nodes pay $G$, and the date-zero operating action retains the original positive-duration sign restriction. A charge for an early forced exit is not introduced.

Let $R_n$ contain all live grid nodes reachable at date $n$ from $x_0$ using any baseline action and either endpoint law, without voluntary surrender. This support set includes paths that no optimal policy uses. It is closed under feasible live transitions and works simultaneously for both regimes and for every $\lambda\in[0,0.25]$.

\begin{proposition}[A statewise enforcement frontier]
\label{prop:r9_enforcement}
For a fixed contract and $m\geq1$, define
\begin{equation}
 F_r^*(m;\lambda,d)=
 \max_{m\leq n<8}\ \max_{x\in R_n\text{ interior}}[G(x)-V^r_n(x;\lambda,d)]_+,
 \label{eq:r9_exact_fee}
\end{equation}
with an empty maximum equal to zero. The voluntary-surrender and noncancellable values coincide at every reachable continuation node if and only if $F\geq F_r^*$. This condition implies equality of both initial class values. It is sufficient, but need not be necessary, for equality only at the initial state.

The threshold is nonincreasing in $m$ and in $d$. Moreover,
\begin{equation}
 F_{\mathrm{adj}}^*(m;\lambda,d)\leq F_0^*(m;\lambda,d).
 \label{eq:r9_capacity_substitute}
\end{equation}
Consequently, deliberate preference adjustment weakly substitutes for enforcement capacity in implementing continued operation, even though its effect on the relative initial-position option can have either sign. If $\underline V_n\leq V^0_n$ uniformly on a parameter region, replacing $V^0$ by $\underline V$ in \eqref{eq:r9_exact_fee} gives a sufficient common charge for both regimes throughout that region.
\end{proposition}

The necessity in this proposition is statewise: an eligible node with $G-F>V_n$ makes equality at that node impossible. The sufficiency follows by backward induction. The order in \eqref{eq:r9_capacity_substitute} follows from inclusion of feasible policies, not from an assumed sign of the four-class option. This distinction matters in the numerical results. A low charge can cause the unadjusted positive class to use surrender disproportionately, reversing the relative option even though adjustment still reduces the compensation required within every class.

The implementation of the regional frontier uses two feasible no-adjustment policies and their all-date Bernstein values. Nonnegative duration makes the lower benefit endpoint $d=0.4$ sufficient. On each of sixteen law subintervals, the minimum restricted coefficient is a policy lower value at each node. Taking the larger of the two policy lower values is still a lower value for $V^0$. Subtracting the arithmetic allowance and maximizing the resulting liquidation deficit over $R_n$ supplies the common charge. This calculation does not require an optimizer at every law probability or a neural policy. It also avoids paying for enforcement at grid nodes that the initial mandate cannot reach.

Free surrender provides the opposite endpoint of this frontier. If $G$ is a supersolution of every feasible one-step operating backup, backward induction gives $W_n=G$ at all surrender-eligible dates when $F=0$. With the first interval compulsory, the agent then purchases only that interval of operation. With surrender additionally available at date zero, the original risk classification is bypassed and the value is simply $G(x_0)$. Neither counterfactual is an implementation of the noncancellable contract.

\subsection{Investment Permissions and Their Shadow Values}
\label{subsec:r9_permissions}
The mandate can separately restrict funded purchases and stock borrowing at the initial date. A long permission $L$ limits risky holdings to $L X_0$; a short permission $S$ limits borrowed risky holdings to $S X_0$. A liquidity-reserve requirement can motivate the former, and a stock-loan or collateral line the latter. Their levels need not be symmetric. The baseline permits $L=0.8$ and $S=0.5$. These numbers are contractual parameters, not estimates of market-wide restrictions. Both $\pi=L$ and $\pi=-S$ carry financial risk.

We vary these permissions only at the first date, leaving all later menus, the $1/8$-year decision horizon, and the transition protocol unchanged. To distinguish a permission effect from a missed action-grid point, we allow every risky share in the relevant interval for each common consumption--adjustment pair. The three original first-date proposals are retained when feasible. Consumption and deliberate adjustment themselves remain on their original finite pairs; the later target is still the complete 1,568-action menu. This is a precisely defined enlargement of initial financial permissions, not continuous optimization of all three controls over the full horizon.

\begin{proposition}[Exact first-date risky-share reduction]
\label{prop:r9_knots}
Fix the optimized continuation $V^r_1$, a common pair $(c,\theta)$, and initial state $x_0=(2,1.25)$. For $\pi\in[-0.6,1]$, every initial proposed increment remains strictly inside the state domain. On this range, the first-date payoff $Q^r(c,\theta,\pi)$ under the finite settlement protocol is continuous and piecewise affine in $\pi$. Its breakpoints are the risky shares at which a conditional arrival wealth coordinate crosses a wealth-grid node. Hence the nonpositive maximum over $[-S,0]$ and the positive supremum over $(0,L]$ are obtained by checking endpoints and these breakpoints for each common $(c,\theta)$ pair, together with feasible frozen proposals. The closure at zero is used only for the positive supremum; a reported positive maximizer must be strictly positive.
\end{proposition}

This reduction replaces a dense risky-share scan by a finite exact search for the specified enlarged initial menu. The numerical implementation contains 121 common consumption--adjustment pairs and 2,689 candidate triples over the wider range. It does not convert finite consumption and adjustment choices into continuous ones. Nor does it convert the piecewise-affine settlement payoff into the diffusion's local mean--variance formula.

On a smooth interpolation cell at a binding optimum, let
\begin{equation}
 \mu_L^r=\partial_\pi Q^r(c_+^*,\theta_+^*,L),\qquad
 \mu_S^r=-\partial_\pi Q^r(c_-^*,\theta_-^*,-S).
 \label{eq:r9_permission_prices}
\end{equation}
These are one-sided values of relaxing the long and short permissions. At a switch between optimizing actions the appropriate directional envelope replaces a single derivative. For initial-only nested permission sets, the positive class value is nondecreasing in $L$, and the nonpositive class value is nondecreasing in $S$, regardless of whether a differentiable multiplier exists.

At a regular indifference locus with $D_A^r=\mathcal A^r_+-\mathcal A^r_-\neq0$, the envelope identities give
\begin{equation}
 \frac{\partial d_r^*}{\partial F}=\frac{\mathcal S^r_+-\mathcal S^r_-}{D_A^r},\qquad
 \frac{\partial d_r^*}{\partial L}=-\frac{\mu_L^r}{D_A^r},\qquad
 \frac{\partial d_r^*}{\partial S}=\frac{\mu_S^r}{D_A^r}.
 \label{eq:r9_frontier_derivatives}
\end{equation}
These are local statements under the exposed-feature and nonzero-slope conditions, not global monotonicity assertions about a difference of convex functions in $d$. They separate three economic margins: duration values an operating benefit, surrender exposure values a termination charge, and the first-date continuation gradient values a binding portfolio permission. Preference adjustment changes these objects through reoptimization. It is not a substitute for observing which permission actually binds.

\subsection{Certification of the Expanded Contract Family}
\label{subsec:r9_certificate}
Adding surrender is compatible with the transition-law certificate. It is an additional action with reward $G-F$ and zero live kernel. It therefore preserves affine dependence on the law probability, nonnegative substochastic kernels, and the finite-horizon information ordering. A fixed policy has the representation
\begin{equation}
 J_p(\lambda,d,F)=\sum_{j=0}^{8} b_{j,8}(\lambda)
       \bigl(B_{p,j}+d A_{p,j}-F S_{p,j}\bigr),
 \label{eq:r9_tensor_policy}
\end{equation}
where $b_{j,8}$ denotes the degree-eight Bernstein basis. Early surrender merely shortens the effective polynomial degree; degree elevation retains the common representation. The upper chord and count recursions use the same expanded menu, including the same date and state eligibility of surrender. Feasible lower policies retain their own surrender choices when evaluated at another fee.

For each law interval, upper values at the four $(d,F)$ corners give a tensor upper bound by convexity of the optimized value in these two reward parameters. Each fixed-policy lower value is affine in them. Restricted Bernstein coefficients then bound the two initial class differences uniformly on a three-parameter box. The maximum over lower policies is taken only after a bound valid for each individual policy has been constructed. The full argument, including arithmetic, is given in Supplement S.13.

There are three separate targets. The mathematical settlement contract uses real arithmetic; the implemented certificate concerns its stored transition and reward arrays; and the diffusion or another settlement contract is a different economic target. Let $\epsilon^r_{\sigma,n}$ be a bound for transferring a class value between two specified finite-array targets. A sufficient backward error account is
\begin{equation}
 \epsilon^r_{\sigma,8}=\epsilon_G,\qquad
 \epsilon^r_{\sigma,n}\leq\epsilon_{r,n}+\epsilon_{K,n}B_{n+1}
                           +\beta_n\epsilon^r_{\sigma,n+1},
 \label{eq:r9_target_error}
\end{equation}
where reward errors include surrender and boundary payments, $\epsilon_K$ is a row-$\ell^1$ kernel error, $B$ bounds continuation magnitude, and $\beta$ bounds row mass. The two class errors must be added to the error in a difference certificate. The reported arithmetic allowance covers evaluation of the stored arrays, not an uncomputed $\epsilon_K$ for the exact constructor or for a diffusion approximation. Selected-control reconstruction checks consistency, but does not estimate an interval enclosure of every transcendental operation. No target transfer is inferred from a small reconstruction discrepancy.
